% ============================================================================
% Patient-Tap paper -- revised sections (drop-in for IEEEtran conference).
%
% Every number below was measured from the patient-tap-android code on
% 2026-10-20 (Flutter 3.44.5, Dart 3.12.2). Re-run `flutter test` before
% submission if the code changes.
%
% Required packages in the preamble (most IEEE templates already load them):
%   \usepackage{amsmath}
%   \usepackage{booktabs}   % optional; swap \toprule etc. for \hline if absent
%   \usepackage{url}
%
% Lines marked TODO need input from the team.
% ============================================================================


% ----------------------------------------------------------------------------
% AUTHOR BLOCK  (replace the existing \author{...})
% ----------------------------------------------------------------------------
\author{
\IEEEauthorblockN{Viral Dhoka}
\IEEEauthorblockA{\textit{Dept. of Computer Engineering}\\
\textit{Vishwakarma Institute of Technology}\\
Pune, India\\
viral.1251070777@vit.edu}
\and
\IEEEauthorblockN{Tanaya Thorat}
\IEEEauthorblockA{\textit{Dept. of Computer Engineering}\\
\textit{Vishwakarma Institute of Technology}\\
Pune, India\\
tanaya.1251070438@vit.edu}
\and
\IEEEauthorblockN{Samarth Dhage}
\IEEEauthorblockA{\textit{Dept. of Computer Engineering}\\
\textit{Vishwakarma Institute of Technology}\\
Pune, India\\
samarth.1251070384@vit.edu}
\and
\IEEEauthorblockN{Vyanktesh Yeole}
\IEEEauthorblockA{\textit{Dept. of Computer Engineering}\\
\textit{Vishwakarma Institute of Technology}\\
Pune, India\\
vyanktesh.1251070836@vit.edu}
\and
\IEEEauthorblockN{Samruddhi Wankhede}
\IEEEauthorblockA{\textit{Dept. of Computer Engineering}\\
\textit{Vishwakarma Institute of Technology}\\
Pune, India\\
samruddhi.1251070439@vit.edu}
\and
\IEEEauthorblockN{Vedant Patil}
\IEEEauthorblockA{\textit{Dept. of Information Technology}\\
\textit{Vishwakarma Institute of Technology}\\
Pune, India\\
vedant.1251130247@vit.edu}
% TODO: add the faculty guide as a seventh author or in the Acknowledgment,
% whichever your department expects.
}


% ----------------------------------------------------------------------------
% REPLACES Section IV-C "Payload Design"
% ----------------------------------------------------------------------------
\subsection{Payload Format}
\label{sec:payload}

The tag carries a single NDEF record of type name format \emph{media} with
the MIME type \texttt{application/x-patient-tap} [6]. Its payload is built in
three layers.

\emph{1) Compact projection.} The application does not serialise the full
profile. It projects the fields the user selected onto a JSON object with
one- and two-letter keys (Table~\ref{tab:fields}). Medications, contacts and
treatment-log entries are stored as positional arrays rather than objects, which
removes repeated key names. Every payload carries a schema version
\texttt{v} (currently 2) so that later formats can be told apart.

\begin{table}[htbp]
\caption{Fields in the Compact Tag Projection}
\label{tab:fields}
\centering
\begin{tabular}{@{}lll@{}}
\toprule
Key & Content & Included when \\
\midrule
\texttt{v}   & Schema version (2)                  & Always \\
\texttt{id}, \texttt{n}, \texttt{bt} & Patient ID, name, blood type & Basic info \\
\texttt{a}   & List of allergies                   & Allergies \\
\texttt{m}   & [name, dose, frequency] per drug    & Medications \\
\texttt{dnr} & Do-not-resuscitate flag (0/1)       & DNR \\
\texttt{c}   & [name, phone] per contact           & Contacts \\
\texttt{l}   & [time, responder ID, action]        & Treatment log \\
\bottomrule
\end{tabular}
\end{table}

\emph{2) Authenticated encryption.} The UTF-8 JSON string $P$ is encrypted
with AES-256-GCM [1] under key $K$ and a fresh 96-bit nonce $N$ drawn from the
platform's cryptographically secure random number generator. No associated
data is used. The result is packed as
\begin{equation}
B = N_{[12]} \,\|\, C_{[|P|]} \,\|\, T_{[16]},
\label{eq:wire}
\end{equation}
so the binary size is $|B| = |P| + 28$ bytes. This layout is byte-for-byte
identical to the output of the browser Web Crypto API, so tags written by the
project's web prototype and its Android application are interchangeable.

\emph{3) Text encoding.} $B$ is Base64-encoded and stored as the NDEF record
payload, which expands it by a factor of $4/3$.

On read, the application first looks for a record with the Patient-Tap MIME
type. This lets it ignore other records that a card may already carry, such
as a manufacturer URL. It then Base64-decodes the payload, rejects anything
shorter than 28 bytes, and attempts GCM decryption. Any authentication-tag
mismatch is reported to the responder as a tampered or unreadable card. No
plaintext is shown unless verification succeeds.

\emph{Capacity.} The target tag is the NXP NTAG215, which offers 504 bytes of
user memory [7]. Its capability container declares a 496-byte NDEF data area,
and the NDEF TLV wrapper takes 4 of those, so the largest NDEF message it
holds is 492\,B. For the built-in demonstration profile with every category
selected, $|P| = 257$\,B, $|B| = 285$\,B, the Base64 payload is 380\,B, and
the complete NDEF record is 411\,B, which fits on the tag. With only basic
information and allergies selected, $|B| = 113$\,B. As the user selects
categories, a live meter shows the size of the NDEF message that would be
written, computed as
\begin{equation}
S(|B|) = h + |\tau| + 4\left\lceil |B|/3 \right\rceil,
\label{eq:ndef}
\end{equation}
where $|\tau| = 25$\,B is the MIME type and the record header $h$ is 3\,B
for a payload of up to 255\,B and 6\,B above that. A unit test checks
Eq.~\eqref{eq:ndef} against the record the NFC library actually builds. The
application refuses to write a message larger than 492\,B, and again refuses
if the message exceeds the size the physical tag reports. Writes are retried
up to three times, because a hand moving mid-write is the most common cause
of failure.

Because of the Base64 expansion and record header, the largest encrypted
payload that fits an NTAG215 is $|B| = 345$\,B, that is 317\,B of JSON,
rather than the full 504\,B. Writing $B$ directly as a binary record would
recover about 115\,B, and we plan to do this in the next schema version.
An earlier build compared $|B|$ itself against 504\,B, so payloads between
346 and 504\,B passed the meter and then failed at the tag. Measuring the
encoded record removed this mismatch.


% ----------------------------------------------------------------------------
% REPLACES Section V-A "Authenticated Encryption" (key paragraph onwards)
% ----------------------------------------------------------------------------
\subsection{Key Management in the Prototype}
\label{sec:key}

AES-GCM protects a payload only as well as its key is protected. We
therefore describe the prototype's key handling exactly as it is implemented.

\emph{Tag key.} To let any installation of the application act as a
responder without an account or network connection, the prototype uses one
\emph{shared responder key} for all tags:
\begin{equation}
K = \mathrm{SHA\text{-}256}(s),
\label{eq:key}
\end{equation}
where $s$ is a fixed string compiled into the application. Every copy of the
app derives the same $K$. This makes the offline, no-sign-in responder
workflow possible. It also means that the encryption provides
\textbf{integrity, not confidentiality, against an attacker who has the
application}: anyone who extracts $s$ from the APK can decrypt any tag and
can also create tags that verify correctly. In this prototype, the
confidentiality guarantee holds only against readers who do not have the
application, such as a generic NFC reader app.

We consider this acceptable for a feasibility prototype and not acceptable
for deployment. A deployable design needs (i) keys stored in hardware-backed
storage such as the Android Keystore [8], (ii) per-patient or per-card keys
rather than one global key, and (iii) a way for authorised responders to
obtain those keys. Section~\ref{sec:future} outlines options, each of which
trades off against the offline goal.

\emph{Share-link keys.} Emergency share links, which are sent through the
optional cloud service, do not use $K$. Each link is encrypted under a fresh
random 256-bit key that travels only in the URL fragment and so never reaches
the server. A unit test confirms that one link's key cannot decrypt another
link (Table~\ref{tab:tests}).

\emph{Nonces.} Nonces are 96 bits drawn independently at random for every
write. Because all tags share one key, the risk of nonce reuse grows with the
total number of writes across all users. Random 96-bit nonces stay within
NIST's recommended bound of $2^{32}$ invocations per key [1], but a per-card
key would remove this shared exposure entirely. Nonce reuse under GCM leaks
the authentication subkey and enables forgeries [9].

\emph{Local PIN.} The application lock stores $\mathrm{SHA\text{-}256}(\mathit{salt}
\,\|\, \texttt{:} \,\|\, \mathit{PIN})$ with a random 128-bit salt and never
stores the PIN itself. The salt defeats precomputed tables, but a single fast
hash over a 4--6 digit PIN can still be brute-forced offline in well under a
second. Production builds should use a memory-hard function such as Argon2id
[10], [11] together with attempt limits.


% ----------------------------------------------------------------------------
% REPLACES Section VII "Evaluation Methodology"
% ----------------------------------------------------------------------------
\section{Evaluation}
\label{sec:eval}

\subsection{Automated Tests}

The security- and safety-critical logic is covered by 20 unit tests in the
project repository [2]. All 20 pass on Flutter~3.44.5 / Dart~3.12.2
(Table~\ref{tab:tests}). The tests exercise the production code directly;
the cryptographic primitives come from the \texttt{cryptography} Dart
package.

\begin{table}[htbp]
\caption{Unit-Test Results (20/20 Passing)}
\label{tab:tests}
\centering
\begin{tabular}{@{}p{0.27\columnwidth}p{0.63\columnwidth}@{}}
\toprule
Area (tests) & Property verified \\
\midrule
Encryption (7) &
Round-trip correctness; output matches Eq.~\eqref{eq:wire} and the size
estimator; a fresh nonce on every call (identical plaintexts give different
ciphertexts); a flipped bit in the tag is rejected; a truncated payload is
rejected; share-link keys are mutually exclusive; the PIN hash is salted,
deterministic and does not contain the PIN. \\
\addlinespace
Tag payload (6) &
All fields survive projection and parsing; a partial selection writes only
the chosen categories; the demo profile fits the NTAG215 budget; the size
estimate of Eq.~\eqref{eq:ndef} equals the NDEF library's record length on
both sides of the short-record boundary; a payload whose binary fits but
whose encoded record does not is rejected; malformed JSON yields an empty
profile instead of a crash. \\
\addlinespace
Conflict checker (7) &
Direct allergy match is flagged as danger; drug-class cross-reactivity
(e.g.\ amoxicillin with a penicillin allergy) is flagged as danger; a repeat
dose within 4\,h triggers a warning, one older than 4\,h does not; an
interaction with a current medication triggers a warning; an unrelated drug
and an empty input pass. \\
\bottomrule
\end{tabular}
\end{table}

\subsection{Payload Size}

Table~\ref{tab:sizes} gives the measured sizes for the demonstration profile
(one patient with two allergies, two medications with dose and frequency, and
two emergency contacts). Encryption adds a constant 28\,B; Base64 and the
NDEF header add the rest.

\begin{table}[htbp]
\caption{Measured Payload Sizes (Demo Profile)}
\label{tab:sizes}
\centering
\begin{tabular}{@{}lrr@{}}
\toprule
Stage & All fields & Basic + allergies \\
\midrule
JSON plaintext $|P|$          & 257\,B & 85\,B  \\
Encrypted binary $|B|$        & 285\,B & 113\,B \\
Base64 payload                & 380\,B & 152\,B \\
NDEF record                   & 411\,B & 180\,B \\
NTAG215 NDEF capacity         & \multicolumn{2}{c}{492\,B} \\
\bottomrule
\end{tabular}
\end{table}

\subsection{Device Testing}

% TODO: fill in from real runs. Suggested minimum so the paper has
% measurements, not just a plan:
%   - Devices: model, Android version (aim for >= 3 phones).
%   - Tags: NTAG213 / 215 / 216, card and wristband form factors.
%   - For each device x tag: 20 write+read cycles in airplane mode; report
%     success rate and median / 95th-percentile time from tap to display.
%   - Tamper check: edit bytes on a written tag with a generic NFC app and
%     confirm the responder screen rejects it.
The end-to-end write, scan and display path was exercised in airplane mode on
\textit{[TODO: devices]} with \textit{[TODO: tag types]}. Over
\textit{[TODO: N]} cycles we observed a success rate of
\textit{[TODO]}\,\% and a median tap-to-display time of \textit{[TODO]}\,ms.


% ----------------------------------------------------------------------------
% ADD to Section IX "Future Scope". Put \label{sec:future} directly after
% that section's \section{...} line so the reference in Sec. V resolves.
% ----------------------------------------------------------------------------
Replacing the shared responder key is the main open problem. Options under
consideration include: per-card keys printed as a QR code on the card itself,
so that the physical card is the credential; keys issued to verified
responders through a hospital or EMS identity service and cached for offline
use; and a split design in which a small unencrypted ``critical'' subset (for
example blood type and life-threatening allergies) is readable by anyone,
while the full record requires an authorised key.


% ----------------------------------------------------------------------------
% REFERENCES  (replaces the whole bibliography)
% ----------------------------------------------------------------------------
\begin{thebibliography}{00}

\bibitem{b1} M. Dworkin, ``Recommendation for block cipher modes of
operation: Galois/Counter Mode (GCM) and GMAC,'' NIST Special Publication
800-38D, National Institute of Standards and Technology, 2007,
doi: 10.6028/NIST.SP.800-38D.

\bibitem{b2} V. Dhoka, T. Thorat, S. Dhage, V. Yeole, S. Wankhede, and
V. Patil, ``Patient-Tap: Source code and test suite,'' GitHub repository.
[Online]. Available: \url{https://github.com/viralala/Patient-Tap}.
Accessed: Oct. 20, 2026.

\bibitem{b3} Google, ``Flutter documentation.'' [Online]. Available:
\url{https://docs.flutter.dev/}. Accessed: Oct. 20, 2026.

\bibitem{b4} Supabase, ``Row Level Security.'' [Online]. Available:
\url{https://supabase.com/docs/guides/database/postgres/row-level-security}.
Accessed: Oct. 20, 2026.

\bibitem{b5} Google, ``ML Kit Text Recognition v2.'' [Online]. Available:
\url{https://developers.google.com/ml-kit/vision/text-recognition/v2}.
Accessed: Oct. 20, 2026.

\bibitem{b6} NFC Forum, ``NFC Data Exchange Format (NDEF) Technical
Specification,'' NFC Forum, Wakefield, MA, USA.

\bibitem{b7} NXP Semiconductors, ``NTAG213/215/216 -- NFC Forum Type~2 Tag
compliant IC with 144/504/888 bytes user memory,'' Product data sheet.

\bibitem{b8} Android Developers, ``Android Keystore system.'' [Online].
Available: \url{https://developer.android.com/privacy-and-security/keystore}.
Accessed: Oct. 20, 2026.

\bibitem{b9} H. B\"{o}ck, A. Zauner, S. Devlin, J. Somorovsky, and
P. Jovanovic, ``Nonce-disrespecting adversaries: Practical forgery attacks on
GCM in TLS,'' in \emph{Proc. 10th USENIX Workshop on Offensive Technologies
(WOOT)}, 2016.

\bibitem{b10} A. Biryukov, D. Dinu, D. Khovratovich, and S. Josefsson,
``Argon2 memory-hard function for password hashing and proof-of-work
applications,'' RFC 9106, Internet Engineering Task Force, Sept. 2021.

\bibitem{b11} OWASP Foundation, ``Password Storage Cheat Sheet.'' [Online].
Available:
\url{https://cheatsheetseries.owasp.org/cheatsheets/Password_Storage_Cheat_Sheet.html}.
Accessed: Oct. 20, 2026.

\bibitem{b12} D. McGrew, ``An interface and algorithms for authenticated
encryption,'' RFC 5116, Internet Engineering Task Force, Jan. 2008.

% TODO: add 5-8 related-work papers on NFC/QR medical IDs, offline mHealth
% records and emergency data access, found through IEEE Xplore or Google
% Scholar, and discuss them in Section III. Do not cite a paper you haven't
% opened and checked.

\end{thebibliography}
